% GENERATED by scripts/build-lecture-slides.py; do not edit. Portable profile: XeTeX, theme and fonts from the core.
\makeatletter\IfFileExists{core/beamer/beamerthememx.sty}{\def\mxcorepath{core/}\def\input@path{{core/beamer/}{core/}}}{\def\mxcorepath{../core/}\def\input@path{{../core/beamer/}{../core/}}}\makeatother
\documentclass[aspectratio=169,10pt,t]{beamer}
\usetheme{mx}
\usepackage{graphicx,booktabs,mathtools,array,adjustbox}
\usepackage[all]{xy}
\providecommand{\actualtextwidth}{\textwidth}
\providecommand{\veccor}[2]{\left(\begin{matrix}#1\\#2\end{matrix}\right)}
\DeclarePairedDelimiter{\VecNorm}{\lVert}{\rVert}
\providecommand{\ETC}{\ldots}
\usepackage[french]{babel}
\setlength{\parskip}{.45em}
\hypersetup{pdfauthor={},pdfsubject={Lecture slides},pdfkeywords={},pdftitle={Arithmétique}}
\providecommand{\N}{\mathbb{N}}
\providecommand{\Z}{\mathbb{Z}}
\providecommand{\D}{\mathbb{D}}
\providecommand{\Q}{\mathbb{Q}}
\providecommand{\R}{\mathbb{R}}
\providecommand{\C}{\mathbb{C}}
\providecommand{\vect}[1]{\overrightarrow{#1}}
\providecommand{\vi}{\vec{\imath}}
\providecommand{\vj}{\vec{\jmath}}
\providecommand{\vk}{\vec{k}}
\providecommand{\norm}[1]{\left\lVert #1\right\rVert}
\providecommand{\vc}[2]{\begin{pmatrix}#1\\#2\end{pmatrix}}
\providecommand{\deter}[2]{\det\left(#1,#2\right)}
\providecommand{\abs}[1]{\left\lvert #1\right\rvert}
\providecommand{\dd}{\,\mathrm{d}}
\providecommand{\e}{\mathrm{e}}
\providecommand{\Card}{\operatorname{Card}}
\providecommand{\E}{\operatorname{E}}
\providecommand{\V}{\operatorname{V}}
\providecommand{\degre}{^{\circ}}
\providecommand{\un}[1]{\,\mathrm{#1}}
\providecommand{\sci}[2]{#1\times10^{#2}}
\providecommand{\pgcd}{\operatorname{PGCD}}
\providecommand{\ppcm}{\operatorname{PPCM}}
\title{Arithmétique}
\subtitle{Cours et exemples}
\date{}
\author{}
\begin{document}
\mxtitleframe
\begin{frame}[allowframebreaks]{Introduction}
Soient \(a\) et \(b\) deux entiers naturels.
On dit que \(a\) divise \(b\) ou que \(a\) est un diviseur de \(b\) s'il existe un entier naturel \(q\) tel que \(b = aq\).
On dit aussi que \(b\) est un multiple de \(a\) ou encore que \(b\) est divisible par \(a\).

L'ensemble des diviseurs de \(a\) est noté \(D_{(a)}\) ou \(D_a\).
\end{frame}
\begin{frame}[allowframebreaks]{Nombres premiers}
\begin{enumerate}
\setcounter{enumi}{0}
                        \item Un nombre \(n\in\mathbb{N}^*\) est dit \emph{premier} 
                              s'il admet exactement deux diviseurs positifs.
                        
\item Deux nombres entiers relatifs \(a\) et \(b\) sont dits \emph{premiers entre eux} s'ils
                              n'ont pas de diviseur commun positif autre que \(1\).

\end{enumerate}
\end{frame}
\begin{frame}[allowframebreaks]{Nombres premiers (2)}
Le nombre \(1\) \textbf{n'}est \textbf{pas} premier,
                        le plus petit nombre premier est \(2\).
\end{frame}
\begin{frame}[allowframebreaks]{Nombres premiers}
\begin{itemize}

                        \item Tout nombre entier naturel \(n > 1\) qui n'est pas premier 
                              admet au moins un diviseur premier \(p \leq \sqrt{n}\).
                        
\item Tout nombre entier naturel \(n \geq 2\) peut se décomposer en produit de facteurs premiers,
                              cette décomposition est unique à l'ordre près.

\end{itemize}
\end{frame}
\begin{frame}[allowframebreaks]{Nombres premiers (2)}
\footnotesize
example
                        Décomposer \(90\) et \(140\) en produits de facteurs premiers.
                        \[
                                    \begin{array}{r|l}
                                           90 & 2 \\
                                           45 & 3 \\
                                           15 & 3 \\
                                            5 & 5 \\
                                            1 &   \\
                                    \end{array}
                                    \qquad
                                    \begin{array}{r|l}
                                          140 & 2 \\
                                           70 & 2 \\
                                           35 & 5 \\
                                            7 & 7 \\
                                            1 &   \\
                                    \end{array}
                        \]
\end{frame}
\begin{frame}[allowframebreaks]{Nombres premiers (3)}
\footnotesize
Ainsi, \(90 = 2 \cdot 3^2 \cdot 5\) et \(140 = 2^2 \cdot 5 \cdot 7\).
                  example
            
            Etant donné un nombre entier naturel \(n > 2\),
            pour savoir si \(n\) est premier,
            il n'est pas nécessaire de vérifier qu'aucun nombre \(q < n\) ne divise \(n\).
            Il suffit de faire la vérification pour les nombres premiers inférieurs ou égaux à \(\sqrt{n}\).
            
                  example
                        Vérifions si \(139\) est premier.
                        
                        Commençons par trouver \(\sqrt{139} \simeq 11,8\).
                        Listons tous les nombres premiers inférieurs à \(11,8\) : 
                              2, 3, 5, 7 et 11.
                        En effectuant la division de \(139\) par ces nombres, on constate qu'aucun d'eux ne divise \(139\).
                        On conclut que \(139\) est premier.
                  example
\end{frame}
\begin{frame}[allowframebreaks]{Nombres premiers}
Il y a une infinité de nombres premiers.
\end{frame}
\begin{frame}[allowframebreaks]{Nombres premiers}
Le crible d'Eratosthenes est une méthode pour trouver tous les nombres premiers jusqu'à un rang \(n\) donné.
            On commence par écrire tous les nombres entiers de \(1\) à \(n\) dans une table 
            contenant \(\lfloor \sqrt{n} \rfloor\)
             \(\lfloor X \rfloor\) ou \(\mathrm{E}\left( X \right)\) signifie 
                       partie entière de \(X\) ;

c'est-à-dire le plus grand entier inférieur ou égal à \(X\).
\end{frame}
\begin{frame}[allowframebreaks]{Nombres premiers}
colonnes.
            On raye le nombre \(1\) puisqu'il n'est pas premier.
            Le prochain nombre non-rayés est premier, on le prend et on raye tous ses multiples.
            On recommence le procédé sur toute la première ligne.
            On s'arrête après avoir traité les nombres premiers inférieurs ou égaux à \(\sqrt{n}\).
            À la fin de l'algorithme, tous les nombres non rayés sont premiers.
\end{frame}
\begin{frame}[allowframebreaks]{Diviseurs et multiples d'un nombre entier relatif}
Soient \(a\) et \(b\) deux nombres entiers relatifs.
On dit que \(a\) divise \(b\), et on note \(a | b\), s'il existe un nombre entier relatif \(q\) tel que \(b = q a\).
On dit aussi que \(b\) est un multiple de \(a\); 
le nombre \(q\) est appelé \emph{quotient} de \(b\) par \(a\).
\end{frame}
\begin{frame}[allowframebreaks]{Diviseurs et multiples d'un nombre entier relatif}
Soient \(a\), \(b\) et \(c\) trois nombres entiers relatifs non-nuls.
\end{frame}
\begin{frame}[allowframebreaks]{Diviseurs et multiples d'un nombre entier relatif (2)}
\begin{enumerate}
\setcounter{enumi}{0}\item Si \(a|b\) et \(b|a\) alors \(a = b\) ou \(a = -b\).
      
\item 
\begin{enumerate}
\setcounter{enumi}{0}
            \item \(a|a\),
            
\item \((a|b \text{ et } b|c) \Rightarrow a|c\),
            
\item \((a|b \text{ et } b|a) \Rightarrow |a| = |b|\).

\end{enumerate}

\end{enumerate}
\end{frame}
\begin{frame}[allowframebreaks]{Diviseurs et multiples d'un nombre entier relatif (3)}
\begin{enumerate}
\setcounter{enumi}{5}\item 
\begin{enumerate}
\setcounter{enumi}{0}
            \item Si \(a|b\) alors \(a|bc\),
            
\item si \(a|b\) alors \(ac|bc\),
            
\item si \(a|b\) et \(a|c\) alors \(a|(\beta b + \gamma c)\), pour tout \((\beta, \gamma) \in \mathbb{Z}^2\).

\end{enumerate}

\item \(1|a\) et \(a|0\).
      
\item \(b|a \text{ et } a\neq 0 \Rightarrow 0 < |b| \leq |a|\).

\end{enumerate}
\end{frame}
\begin{frame}[allowframebreaks]{Diviseurs et multiples d'un nombre entier relatif}
Soit \((a,b) \in \mathbb{Z} \times \mathbb{Z}^*\).

Il existe alors un unique couple d'entiers relatifs \((q,r)\) tel que

\( a = b q + r \text{ et } 0 \leq r < |b|      \;. \)

\begin{itemize}

      \item \(q\)  est le quotient de la division de \(a\) par \(b\),      
      
\item \(r\)  est le reste de la division de \(a\) par \(b\).      

\end{itemize}
\end{frame}
\begin{frame}[allowframebreaks]{Diviseurs et multiples d'un nombre entier relatif (2)}
\(b|a\) si et seulement si \(r=0\).
      
      example
      Prenons \(a = -17\) et \(b = -4\); on a \(q = 5\) et \(r = 3\).
            On vérifie que \(-17 = (-4) \cdot 5 + 3\).
      example
\end{frame}
\begin{frame}[allowframebreaks]{Numération à base \(b\)}
Rappelons que dans l'écriture usuelle d'un nombre entier, on utilise la base \(10\).
      On se sert des symboles : \(0\), \(1\), \(2\), \(3\), \(4\), \(5\), \(6\), \(7\), \(8\) et \(9\) (il y a dix symboles).
\end{frame}
\begin{frame}[allowframebreaks]{Numération à base \(b\) (2)}
Soit \(N\) un entier naturel et
            \( a_n a_{n-1} \dots a_2 a_1 a_0 \)
      où \(a_i \in \{ 0, 1, 2, \dots , 9 \}\) est son écriture décimale (c'est-à-dire usuelle).
      Cette écriture signifie que 
            \[\begin{aligned}
            N & = a_{n} \cdot 10^{n} + a_{n-1} \cdot 10^{n-1} + \cdots 
                  + a_{2} \cdot 10^{2} + a_{1} \cdot 10^{1} + a_{0} \cdot 10^{0}            \\
              & = \sum_{i=0}^{n} a_{i} \cdot 10^{i}
            \end{aligned}\]
\end{frame}
\begin{frame}[allowframebreaks]{Numération à base \(b\)}
Considérons le nombre \(N = 9\;235\;819\).
            On a alors 
            \(N = 9\cdot 10^{6} + 2\cdot 10^{5} + 3\cdot 10^{4} + 5\cdot 10^{3} + 8\cdot 10^{2} + 1\cdot 10^{1} + 9\cdot 10^{0}\).
\end{frame}
\begin{frame}[allowframebreaks]{Numération à base \(b\)}
Pour écrire un nombre entier naturel \(N\) en base \(b\) (où \(b\geq 2\)), on se sert de \(b\) symboles.
      Si l'écriture en base \(b\) de \(N\) est 
            \( a_n a_{n-1} \dots a_2 a_1 a_0 \)
      alors 
            \[\begin{aligned}
            N & = a_{n} \cdot b^{n} + a_{n-1} \cdot b^{n-1} + \cdots 
                  + a_{2} \cdot b^{2} + a_{1} \cdot b^{1} + a_{0} \cdot b^{0}            \\
              & = \sum_{i=0}^{n} a_{i} \cdot b^{i}                                       .
            \end{aligned}\]
\end{frame}
\begin{frame}[allowframebreaks]{Numération à base \(b\) (2)}
Pour indiquer que l'écriture est en base \(b\), on note \( \overline{a_n a_{n-1} \dots a_2 a_1 a_0}^{(b)}\).
      conversionexample
      Écrivons les nombres 13 et 41 en base 2 (écriture binaire).

Ainsi, \(13 = \overline{1101}^{_{(2)}}\).
            
            Ainsi, \(41 = \overline{101001}^{_{(2)}}\).
      conversionexample
\end{frame}
\begin{frame}[allowframebreaks]{Numération à base \(b\)}
conversionexample
            Écrivons les nombres 957 et 188899 en base 8 :

            Ainsi, \(957 = \overline{1675}^{_{(8)}}\).
            
            Ainsi, \(188899 = \overline{560743}^{_{(8)}}\).
      conversionexample
\end{frame}
\begin{frame}[allowframebreaks]{Numération à base \(b\) (2)}
conversionexample
            Écrivons les nombres 18 647 et 1 304 748 en base 12, pour celà, 
            on aura besoin d'ajouter les symboles \(A\) et \(B\) pour les <<chiffres>> dix et onze respectivement.

            Ainsi, \(18\,647 = \overline{A95B}^{_{(12)}}\).
            
            Ainsi, \(1\,304\,748 = \overline{52B090}^{_{(12)}}\).
      conversionexample
\end{frame}
\begin{frame}[allowframebreaks]{Numération à base \(b\)}
\footnotesize
Il existe d'autres notations pour les nombres en base autres que dix.
            Le nombre \(42\) peut par exemple être noté
                  \( \overline{101010}^{_{(2)}} \),
                  \( \overline{101010}^{_{2}} \),
                  \( (\overline{101010})_{2} \),
                  \( {101010}_{2} \),
                  \ldots{}.
            On choisira, dans ce cours, d'utiliser toujours la notation \( \overline{101010}^{_{(2)}} \).
            
            Pour les bases supérieures à dix,
            on ajoute souvent des lettres de l'alphabet latin dans l'ordre.
            On peut utiliser des lettres minuscules aussi bien que capitales.
\end{frame}
\begin{frame}[allowframebreaks]{Congruence}
Soient \((x,y) \in \mathbb{Z}^2\) et \(n\in\mathbb{N}\).
            On dit que \(x\) est congru à \(y\) modulo \(n\) si
            \(n\) divise \(x-y\). 
            On note \(x \equiv y \quad \mathrm{mod}\ n\) ou \(x \equiv y \ [n]\).
            
                  \[\begin{array}{rcl}
                        x \equiv y \quad \mathrm{mod}\ n & \iff & n | x - y                                     \\
                                         & \iff & x - y = k \cdot n \quad (k\in\mathbb{Z})      \\
                                         & \iff & x = y + k \cdot n \quad (k\in\mathbb{Z})      
                  \end{array}\]

La congruence modulo \(0\) est l'égalité.
\end{frame}
\begin{frame}[allowframebreaks]{Congruence}
Soit \(n\) un entier naturel.
\end{frame}
\begin{frame}[allowframebreaks]{Congruence (2)}
\begin{enumerate}
\setcounter{enumi}{0}\item La congruence modulo \(n\) est une relation d'équivalence, 
                        c'est-à-dire que, pour tous entiers relatifs \(x\), \(y\) et \(z\), on a :

\item \textbf{Reflexivité :}  
                                    \(x \equiv x \quad \mathrm{mod}\ n\),

\end{enumerate}
\end{frame}
\begin{frame}[allowframebreaks]{Congruence (3)}
\begin{enumerate}
\setcounter{enumi}{2}\item \textbf{Symétrie :}  
                                    \(x \equiv y \quad \mathrm{mod}\ n \iff y \equiv x \quad \mathrm{mod}\ n\),
                              
\item \textbf{Transitivité :}  
                                    \(\left( x \equiv y \quad \mathrm{mod}\ n \text{ et } y \equiv z \quad \mathrm{mod}\ n \right) \Rightarrow x \equiv z \quad \mathrm{mod}\ n\).

\end{enumerate}
\end{frame}
\begin{frame}[allowframebreaks]{Congruence (4)}
\footnotesize
\begin{enumerate}
\setcounter{enumi}{4}\item La congruence modulo \(n\) est compatible avec l'addition et la multiplication, 
                        c'est-à-dire que, pour tous entiers relatifs \(x_1\), \(x_2\), \(y_1\) et \(y_2\), on a :

                        \[
                              \left.
                                    \begin{array}{l}
                                          x_1 \equiv y_1 \quad \mathrm{mod}\ n \\
                                          x_2 \equiv y_2 \quad \mathrm{mod}\ n
                                    \end{array}
                              \right\rbrace
                              \Rightarrow
                              \left\lbrace
                                    \begin{array}{l}
                                          x_1 + x_2 \equiv y_1 + y_2 \quad \mathrm{mod}\ n             \\
                                          x_1 \cdot x_2 \equiv y_1 \cdot y_2 \quad \mathrm{mod}\ n     \\
                                    \end{array}
                              \right.
                        \]
                        et par conséquent, pour tout \(k\in\mathbb{N}^*\), on a :
                              \(x \equiv y \quad \mathrm{mod}\ n \Rightarrow x^k \equiv y^k \quad \mathrm{mod}\ n\).

\end{enumerate}
\end{frame}
\begin{frame}[allowframebreaks]{Congruence}
\begin{enumerate}
\setcounter{enumi}{0}
                  \item Si \(d|n\) et \(x \equiv y \quad \mathrm{mod}\ n\), alors \(x \equiv y \quad \mathrm{mod}\ d\).
                  
\item Soit \(n\in\mathbb{N}^*\).
                        Pour tout \(x\in\mathbb{Z}\), il existe un unique \(r\in\mathbb{N}\) tel que
                        \(0\leq r < n\) et \(x \equiv r \quad \mathrm{mod}\ n\).

\end{enumerate}
\end{frame}
\begin{frame}[allowframebreaks]{Congruence}
On veut déterminer le reste de la division euclidienne de \(7^{100}\) par \(5\).
            
            On a \(7^{100} = (7^2)^{50} = 49^{50}\).
            Comme \(49 = 50 - 1 = 5 \cdot 10 - 1\),
            alors  \(49 \equiv -1 \quad \mathrm{mod}\ 5\) ; c'est-à-dire \(7^2 \equiv -1 \quad \mathrm{mod}\ 5\).
            Il vient que \((7^2)^{50} \equiv (-1)^{50} \quad \mathrm{mod}\ 5\) et donc \(7^{100} \equiv 1 \quad \mathrm{mod}\ 5\).
            Ainsi, le reste de la division euclidienne de \(7^{100}\) par \(5\) est \(1\).
\end{frame}
\begin{frame}[allowframebreaks]{Congruence}
\footnotesize
On veut déterminer un critère de divisibilité par \(9\).
            
            Soit \(N\) un nombre entier naturel et soit \(a_n a_{n-1} \ldots a_2 a_1 a_0\) son écriture en base 10.
            On a alors \(N = a_n \cdot 10^n + a_{n-1} \cdot 10^{n-1} + \cdots + a_1 \cdot 10 + a_0\).
            
            Montrons, d'abord, que pour tout \(p\in\mathbb{N}\), on a \(10^p \equiv 1 \quad \mathrm{mod}\ 9\) :
                  \[\begin{aligned}
                        10 & = 9 \cdot 1 + 1     \\
                        10 &  \equiv 1 \quad \mathrm{mod}\ 9     \\
                        10^p &  \equiv 1^p \quad \mathrm{mod}\ 9 \\
                        10^p &  \equiv 1 \quad \mathrm{mod}\ 9   \;.
                  \end{aligned}\]
\end{frame}
\begin{frame}[allowframebreaks]{Congruence (2)}
Ainsi, pour tout \(p\in\mathbb{N}\), on a :
                  \[\begin{aligned}
                        10^p &  \equiv 1 \quad \mathrm{mod}\ 9                   \\
                        a_p \cdot 10^p &  \equiv a_p \quad \mathrm{mod}\ 9       \;.
                  \end{aligned}\]

En prenant la somme pour \(p\) allant de \(0\) à \(n\), on a :
                  \[\begin{aligned}
                        a_n10^n+a_{n-1}10^{n-1}+\cdots+a_110+a_0 \equiv a_n+a_{n-1}+\cdots+a_1+a_0 \quad \mathrm{mod}\ 9      \;.
                  \end{aligned}\]
\end{frame}
\begin{frame}[allowframebreaks]{Congruence (3)}
Ainsi, \(N \equiv \sum_{p=0}^n a_p \quad \mathrm{mod}\ 9\).
            Comme \(9|N \iff N \equiv 0 \quad \mathrm{mod}\ 9\), alors \(9|N \iff \sum_{p=0}^n a_p \equiv 0 \quad \mathrm{mod}\ 9\).
            Ainsi, \(N\) est un multiple de \(9\) si et seulement si la somme de ces chiffres est un multiple de \(9\).
\end{frame}
\begin{frame}[allowframebreaks]{Anneau \(\mathbb{Z}/n\mathbb{Z}\)}
Soient \(n\in\mathbb{N}\), \(n \geq 2\), et \(r\in\mathbb{Z}\).
      On définit la classe d'équivalence de \(r\) comme étant 
      l'ensemble des \(x\in\mathbb{Z}\) tel que \(x \equiv r \quad \mathrm{mod}\ n\), 
      on la note par \(\overline{r}\) ou \(\dot{r}\).
      Ainsi,
      \(\overline{r} = \left\lbrace x\in\mathbb{Z} \mid x \equiv r \quad \mathrm{mod}\ n \right\rbrace\).

Bien qu'une classe d'équivalence peut avoir plusieurs représentants, 
      on préfère prendre \( r \in \mathopen{[\![} 0 ; n-1 \mathclose{]\!]}\).
\end{frame}
\begin{frame}[allowframebreaks]{Anneau \(\mathbb{Z}/n\mathbb{Z}\) (2)}
On note par \(\mathbb{Z}/n\mathbb{Z}\) l'ensemble 
      \(\left\lbrace \overline{0},\ \overline{1},\ \ldots ,\ \overline{n-1} \right\rbrace\).
\end{frame}
\begin{frame}[allowframebreaks]{Anneau \(\mathbb{Z}/n\mathbb{Z}\)}
On définit l'addition et la multiplication dans \(\mathbb{Z}/n\mathbb{Z}\) par :       
      \[\begin{aligned}
            \overline{a} + \overline{b} & = \overline{a+b}       & \overline{a} \cdot \overline{b} = \overline{a\cdot b} \;.
      \end{aligned}\]

Ces opérations ne dépendent pas des représentants choisis (compatibilité de la congruence avec l'addition et la multiplication).
      Muni de ces deux opérations, \(\mathbb{Z}/n\mathbb{Z}\) a une structure d'anneau,
\end{frame}
\begin{frame}[allowframebreaks]{Anneau \(\mathbb{Z}/n\mathbb{Z}\)}
c'est-à-dire que pour tout \(\overline{a}\), \(\overline{b}\) et \(\overline{c}\) de \(\mathbb{Z}/n\mathbb{Z}\), on a :
\end{frame}
\begin{frame}[allowframebreaks]{Anneau \(\mathbb{Z}/n\mathbb{Z}\) (2)}
\begin{itemize}
\item \(\overline{a} + \overline{b}\) et \(\overline{a} \cdot \overline{b}\) sont bien définie dans \(\mathbb{Z}/n\mathbb{Z}\),
                  
\item \((\overline{a} + \overline{b}) + \overline{c} = \overline{a} + (\overline{b} + \overline{c})\),
                  
\item \(\overline{a} + \overline{b} = \overline{b} + \overline{a}\),
                  
\item \(\overline{a} + \overline{0} = \overline{a}\),

\end{itemize}
\end{frame}
\begin{frame}[allowframebreaks]{Anneau \(\mathbb{Z}/n\mathbb{Z}\) (3)}
\begin{itemize}
\item il existe un élément, noté \(-\overline{a}\) tel que \(\overline{a} + (-\overline{a}) = \overline{0}\),
                  
\item \((\overline{a} \cdot \overline{b}) \cdot \overline{c} = \overline{a} \cdot (\overline{b} \cdot \overline{c})\),
                  
\item \(\overline{a} \cdot (\overline{b} + \overline{c}) = \overline{a} \cdot \overline{b} + \overline{a} \cdot \overline{c}\),

\end{itemize}
\end{frame}
\begin{frame}[allowframebreaks]{Anneau \(\mathbb{Z}/n\mathbb{Z}\) (4)}
\begin{itemize}
\item \((\overline{b} + \overline{c}) \cdot \overline{a} = \overline{b} \cdot \overline{a} + \overline{c} \cdot \overline{a}\).

\end{itemize}
\end{frame}
\begin{frame}[allowframebreaks]{Anneau \(\mathbb{Z}/n\mathbb{Z}\)}
\(\mathbb{Z}/5\mathbb{Z} = \left\lbrace \overline{0},\ \overline{1},\ \overline{2},\ \overline{3},\ \overline{4} \right\rbrace\).
            Dans ce cas, 
            \[\begin{aligned}
                  \overline{0} & = \left\lbrace \ldots ,\ -15,\ -10,\ -5,\ 0,\ 5,\ 10,\ 15,\ \ldots \right\rbrace ,\\
                  \overline{1} & = \left\lbrace \ldots ,\ -14,\ -9,\ -4,\ 1,\ 6,\ 11,\ 16,\ \ldots \right\rbrace  ,\\
                  \overline{2} & = \left\lbrace \ldots ,\ -13,\ -8,\ -3,\ 2,\ 7,\ 12,\ 17,\ \ldots \right\rbrace  ,\\
                  \overline{3} & = \left\lbrace \ldots ,\ -12,\ -7,\ -2,\ 3,\ 8,\ 13,\ 18,\ \ldots \right\rbrace  ,\\
                  \overline{4} & = \left\lbrace \ldots ,\ -11,\ -6,\ -1,\ 4,\ 9,\ 14,\ 19,\ \ldots \right\rbrace  .
            \end{aligned}\]
\end{frame}
\begin{frame}[allowframebreaks]{Anneau \(\mathbb{Z}/n\mathbb{Z}\) (2)}
\footnotesize
On a les tables d'additions et de multiplications suivantes :
                  \[
                  \begin{array}{|c|c|c|c|c|c|}
                  \hline 
                  + & \overline{0} & \overline{1} & \overline{2} & \overline{3} & \overline{4} \\ 
                  \hline 
                  \overline{0} & \overline{0} & \overline{1} & \overline{2} & \overline{3} & \overline{4} \\ 
                  \hline 
                  \overline{1} & \overline{1} & \overline{2} & \overline{3} & \overline{4} & \overline{0} \\ 
                  \hline 
                  \overline{2} & \overline{2} & \overline{3} & \overline{4} & \overline{0} & \overline{1} \\ 
                  \hline 
                  \overline{3} & \overline{3} & \overline{4} & \overline{0} & \overline{1} & \overline{2} \\ 
                  \hline 
                  \overline{4} & \overline{4} & \overline{0} & \overline{1} & \overline{2} & \overline{3} \\ 
                  \hline 
                  \end{array}
                  \qquad
                  \begin{array}{|c|c|c|c|c|c|}
                  \hline 
                  \times & \overline{0} & \overline{1} & \overline{2} & \overline{3} & \overline{4} \\ 
                  \hline 
                  \overline{0} & \overline{0} & \overline{0} & \overline{0} & \overline{0} & \overline{0} \\ 
                  \hline 
                  \overline{1} & \overline{0} & \overline{1} & \overline{2} & \overline{3} & \overline{4} \\ 
                  \hline 
                  \overline{2} & \overline{0} & \overline{2} & \overline{4} & \overline{1} & \overline{3} \\ 
                  \hline 
                  \overline{3} & \overline{0} & \overline{3} & \overline{1} & \overline{4} & \overline{2} \\ 
                  \hline 
                  \overline{4} & \overline{0} & \overline{4} & \overline{3} & \overline{2} & \overline{1} \\ 
                  \hline 
                  \end{array}
                  .
                  \]
\end{frame}
\begin{frame}[allowframebreaks]{Anneau \(\mathbb{Z}/n\mathbb{Z}\)}
Au lieu de \(\mathbb{Z}/n\mathbb{Z}\), on utilise parfois \(\mathbb{Z}_n\).
\end{frame}
\begin{frame}[allowframebreaks]{Anneau \(\mathbb{Z}/n\mathbb{Z}\)}
\footnotesize
Résolvons dans \(\mathbb{Z}/5\mathbb{Z}\) l'équation \(x+\overline{3} = \overline{2}\).
            Pour celà, additionnons \(\overline{2}\) dans chaque membre de l'égalité, on a alors 
                  \(x + \overline{3} + \overline{2} = \overline{2} + \overline{2} \Rightarrow x = \overline{4}\).
            On aurait aussi pu utiliser la méthode usuelle et trouver \( x = \overline{-1} = \overline{4}\).
            
            Résolvons dans \(\mathbb{Z}/5\mathbb{Z}\) l'équation \(\overline{3}\, x+\overline{3} = \overline{2}\).
            On sait déjà (en utilisant ce qui précède) que \(\overline{3} x = \overline{4}\).
            En multipliant l'égalité par \(\overline{2}\), on obtient 
                  \(\overline{2} \cdot \overline{3} \cdot x = \overline{2} \cdot \overline{4}\)
                  et donc \(x = \overline{3}\).
            
            Résolvons dans \(\mathbb{Z}/5\mathbb{Z} \times \mathbb{Z}/5\mathbb{Z}\) le système d'équations :
            \[
            \left\lbrace
                  \begin{array}{l}
                        \overline{3}\, x + \overline{2}\, y = \overline{1} \\
                        \overline{2}\, x + \overline{4}\, y = \overline{0} \;.
                  \end{array}
            \right.
            \]
\end{frame}
\begin{frame}[allowframebreaks]{Anneau \(\mathbb{Z}/n\mathbb{Z}\) (2)}
On peut additionner les deux équations pour avoir 
                \( ( \overline{3} + \overline{2} ) \, x + ( \overline{2} + \overline{4} ) \, y = \overline{1} + \overline{0} \).
                Ainsi, \(y = \overline{1}\).
            Pour trouver la valeur de \(x\), on peut commencer par multiplier la première équation par \(\overline{3}\),
            \[
            \left\lbrace
                  \begin{array}{l}
                        \overline{4}\, x + \overline{1}\, y = \overline{3} \\
                        \overline{2}\, x + \overline{4}\, y = \overline{0} \\
                  \end{array}
            \right.
            \]
\end{frame}
\begin{frame}[allowframebreaks]{Anneau \(\mathbb{Z}/n\mathbb{Z}\) (3)}
et ensuite additionner les deux équations pour avoir 
                  \( x = \overline{3}\).
\end{frame}
\begin{frame}[allowframebreaks]{PPCM et PGCD}
\footnotesize
Soient \(a\) et \(b\) deux entiers naturels non nuls.
      On appelle \textbf{P}lus \textbf{P}etit \textbf{C}ommun \textbf{M}ultiple 
      ,  en abrégé \emph{PPCM} ,  de \(a\) et \(b\) le nombre 
            \(PPCM(a;b) := \min \left( a\mathbb{N}^* \cap b\mathbb{N}^* \right)\)
            où \(x\mathbb{N}^* := \left\lbrace x \cdot n \mid n \in \mathbb{N}^* \right\rbrace\).

                  Le PPCM de \(a\) et \(b\) est aussi parfois noté \(a \vee b\) ou \([a;b]\).

      Si \(a\) et \(b\) sont deux entiers relatifs, 
      on peut définir le PPCM de \(a\) et \(b\) comme étant le PPCM de \(|a|\) et \(|b|\).
      Si \(a\) et \(b\) sont deux entiers relatifs non-nuls, on a aussi :
            \(PPCM(a;b) := \min \left( \left\lbrace m\in\mathbb{N}^* \mid a|m \text{ et } b|m \right\rbrace \right)\).
\end{frame}
\begin{frame}[allowframebreaks]{PPCM et PGCD}
Soient \(a\) et \(b\) deux entiers naturels non-nuls.

\begin{itemize}

                  \item Si \(a|b\) alors \(PPCM(a;b) = b\),
                  
\item si \(k\in\mathbb{N}^*\) alors \(PPCM(ka;kb) = k \cdot PPCM(a;b)\),
                  
\item si \(c\in\mathbb{N}^*\), \(PPCM(a;b;c) = PPCM(PPCM(a;b);c) = PPCM(a;PPCM(b;c))\)
                        (on peut étendre à un nombre arbitraire d'éléments).

\end{itemize}
\end{frame}
\begin{frame}[allowframebreaks]{PPCM et PGCD}
Pour trouver le PPCM de deux ,  ou de plusieurs ,  nombres, 
            on peut décomposer chacun des nombres en produits de facteurs premiers.

Il suffit ensuite de prendre, pour chaque facteur premier présent, la plus grande puissance.
\end{frame}
\begin{frame}[allowframebreaks]{PPCM et PGCD}
Cherchons \(PPCM(28;24;36)\).
                  
                  Comme 
                  \(
                  \left\lbrace
                  \begin{array}{l}
                        28 = 2^2 \cdot 7 \\
                        24 = 2^3 \cdot 3 \\
                        36 = 2^2 \cdot 3^2 \\
                  \end{array}
                  \right.
                  \)
                  le PPCM de ces trois nombres est \(2^3 \cdot 3^2 \cdot 7 = 504\).
\end{frame}
\begin{frame}[allowframebreaks]{PPCM et PGCD}
Soient \(a\) et \(b\) deux entiers naturels non tous les deux nuls.
      On appelle \textbf{P}lus \textbf{G}rand \textbf{C}ommun \textbf{D}iviseur
      ,  en abrégé \emph{PGCD} ,  de \(a\) et \(b\) le nombre 
            \(PGCD(a,b) := \max \left( D_a \cap D_b \right)\).

                  Le PGCD de \(a\) et \(b\) est aussi parfois noté \(PGCD(a,b)\), \(a \wedge b\) ou \((a;b)\). 
                  On trouve parfois l'acronyme PGDC 
                  (\textbf{P}lus \textbf{G}rand \textbf{D}iviseur \textbf{C}ommun)
\end{frame}
\begin{frame}[allowframebreaks]{PPCM et PGCD (2)}
mais la version la plus répandue est PGCD.

      Deux nombres \(a\) et \(b\) sont dits \emph{premiers entre eux} si \(PGCD(a,b) = 1\).
\end{frame}
\begin{frame}[allowframebreaks]{PPCM et PGCD}
Tout entier divise \(0\), donc \(D_0 \cap D_0\) n'a pas de plus grand élément ;
            on pose, par convention, \(PGCD(0,0)=0\).
            Pour \(a \neq 0\), on a \(PGCD(a,0) = a\).
            Pour des entiers relatifs, on pose \(PGCD(a,b) = PGCD(|a|,|b|)\).
\end{frame}
\begin{frame}[allowframebreaks]{PPCM et PGCD}
Soient \(a\) et \(b\) deux entiers naturels non-nuls.

\begin{itemize}

                  \item Si \(a|b\) alors \(PGCD(a,b) = a\),
                  
\item si \(k\in\mathbb{N}^*\) alors \(PGCD(ka,kb) = k \cdot PGCD(a,b)\),
                  
\item si \(c\in\mathbb{N}^*\), \(\mathrm{PGCD}(a,b,c) = PGCD(PGCD(a,b),c) = PGCD(a,PGCD(b,c))\)
                        (on peut étendre à un nombre arbitraire d'éléments).

\end{itemize}
\end{frame}
\begin{frame}[allowframebreaks]{PPCM et PGCD}
Soient \(a\) et \(b\) deux nombres entiers naturels non-nuls.
            Soit \(r\) le reste de la division de \(b\) par \(a\).
            
            Alors \(PGCD(a,b)\) divise \(r\).
            
            En fait, on a un résultat plus fort : \(PGCD(a,b) = PGCD(a,r)\)
                       (car \(r = b - qa\) et \(b = qa + r\) : les diviseurs communs de \(a\) et \(b\) sont exactement ceux de \(a\) et \(r\)).
\end{frame}
\begin{frame}[allowframebreaks]{PPCM et PGCD}
Soient \(a\) et \(b\) deux entiers naturels non-nuls.
            On a
            \[
                  PPCM(a;b)\cdot PGCD(a,b)=a\cdot b.
            \]
\end{frame}
\begin{frame}[allowframebreaks]{PPCM et PGCD}
Pour trouver le PGCD de deux ,  ou de plusieurs ,  nombres, 
            on peut décomposer chacun des nombres en produits de facteurs premiers.

Il suffit ensuite de prendre, pour chaque facteur premier présent, la plus petite puissance
            (éventuellement \(0\)).
\end{frame}
\begin{frame}[allowframebreaks]{PPCM et PGCD}
Cherchons \(PGCD(36,56)\).
                  
                  Comme 
                  \(
                  \left\lbrace
                  \begin{array}{l}
                        36 = 2^2 \cdot 3^2 \\
                        56 = 2^3 \cdot 7 \\
                  \end{array}
                  \right.
                  \)
                  le PGCD de ces deux nombres est \(2^2 = 4\).
\end{frame}
\begin{frame}[allowframebreaks]{PPCM et PGCD}
Il existe un algorithme ,  celui d'Euclide ,  permettant de trouver le PGCD de deux nombres entiers naturels.

Soient \(a\) et \(b\) deux entiers naturels non-nuls tels que \(a>b\).
            On pose \(a_0=a\), \(a_1=b\) et, pour tout entier naturel \(n\),
            \[
              \begin{aligned}
                a_n &= q_n a_{n+1}+r_n,
                & 0 &\leq r_n<a_{n+1}, \\
                a_{n+2} &= r_n.
              \end{aligned}
            \]
\end{frame}
\begin{frame}[allowframebreaks]{PPCM et PGCD (2)}
Les nombres \(q_n\) et \(r_n\) sont respectivement le quotient et le
            reste de la division de \(a_n\) par \(a_{n+1}\).
\end{frame}
\begin{frame}[allowframebreaks]{PPCM et PGCD}
On peut dresser le tableau suivant
            \[
            \begin{array}{|c|c|c|c|c|}
            \hline 
              & q_0 & q_1 & q_2 & q_3 \\ 
            \hline 
            a_0 = a & a_1 = b & a_2 = r_0 & a_3 = r_1  & \cdots \\ 
            \hline 
            r_0 = r & r_1 = a_1 - a_2 \cdot q_1 & r_2 = a_2 - a_3 \cdot q_2 & \cdots & \cdots \\ 
            \hline 
            \end{array} 
            \]

et le dernier reste (c'est dans la dernière ligne) non-nul est le PGCD de \(a\) et \(b\).
\end{frame}
\begin{frame}[allowframebreaks]{PPCM et PGCD}
Cherchons le PGCD de \(36\) et \(200\).
                  \[
                  \begin{array}{|c|c|c|c|c|}
                  \hline 
                           & q_0 = 5  & q_1 = 1   & q_2 = 1 & q_3 = 4 \\ 
                  \hline 
                  200      &      36  &      20   &      16 &       4 \\ 
                  \hline 
                  r_0 = 20 & r_1 = 16 & r_2 = 4   &       0 &         \\ 
                  \hline 
                  \end{array} 
                  \]

Le dernier reste non-nul est \(4\), donc \(PGCD(36,200) = 4\).
\end{frame}
\begin{frame}[allowframebreaks]{PPCM et PGCD}
Soient \(a\) et \(b\) deux entiers relatifs non-nuls et notons par \(d\) le PGCD de \(a\) et \(b\).
                  On a alors :

\begin{itemize}

                        \item Il existe deux entiers \(u\) et \(v\) vérifiant : \(a u + b v = d\),
                        
\item pour tous entiers \(u\) et \(v\), le nombre \(a u + b v\) est un multiple de \(d\).

\end{itemize}
\end{frame}
\begin{frame}[allowframebreaks]{PPCM et PGCD (2)}
Les nombres \(u\) et \(v\) sont appelés \emph{coefficients de Bézout},
                  ils ne sont pas uniques.
\end{frame}
\begin{frame}[allowframebreaks]{PPCM et PGCD}
On sait que \(PGCD(200,36) = 4\).
                  Cherchons deux entiers relatifs \(u\) et \(v\) tels que \(200 u + 36 v = 4\).
                  
                  \[
                        \begin{array}{|r@{\ =\ }l|r@{\ =\ }l|r@{\ =\ }l|}
                              \hline
                              \multicolumn{2}{|c|}{\text{Division}} &
	                                     \multicolumn{2}{|c|}{\text{Écriture numérique} } & 
	                                           \multicolumn{2}{|c|}{\text{Écriture littérale} } \\
                              \hline
                              200 & 5 \cdot 36 + 20 & 20 & 200 - 5 \cdot 36       & r_0  & a   - 5 \cdot b    \\
                              \hline
                              36  & 1 \cdot 20 + 16 & 16 & 36  - 1 \cdot 20       & r_1  & b   - 1 \cdot r_0  \\
                              \hline
                              20  & 1 \cdot 16 + 4  & 4  & 20  - 1 \cdot 16       & r_2  & r_0 - 1 \cdot r_1  \\
                              \hline
                              16  & 4 \cdot 4  + 0  & \multicolumn{2}{|c|}{}      & \multicolumn{2}{|c|}{}  \\
                              \hline
                        \end{array}
                  \]
\end{frame}
\begin{frame}[allowframebreaks]{PPCM et PGCD (2)}
Ainsi, on a 
                        \[\begin{aligned}
                              r_2 & = (a-5b) - (b-r_0)       \\
                                  & = (a-5b) - (b-(a-5b))    \\
                                  & = a-5b-b+a-5b            \\
                                  & = 2a - 11 b              \;.
                        \end{aligned}\]

D'où \(4 = 2 \cdot 200 - 11 \cdot 36\).
\end{frame}
\begin{frame}[allowframebreaks]{PPCM et PGCD}
Il existe une version “\emph{étendue}” de l'algorithme d'Euclide.
		Elle permet, en même temps que de connaître le PGCD de deux nombres \(a\) et \(b\), 
		de trouver les coefficients de Bézout \(u\) et \(v\).
\end{frame}
\begin{frame}[allowframebreaks]{PPCM et PGCD}
Il est identique au précédent, il faut juste ajouter deux lignes.
		\[
			\begin{array}{|c|c|c|c|c|}
				\hline 
					& q_0 & q_1 & q_2 & \cdots \\ 
				\hline 
					a_0 = a & a_1 = b & a_2 = r_0 & a_3 = r_1  & \cdots \\ 
				\hline 
					r_0 = r & r_1 = a_1 - a_2 \cdot q_1 & r_2 = a_2 - a_3 \cdot q_2 & r_3 = a_3 - a_4 \cdot q_3 & \cdots \\ 
				\hline
					\boldsymbol{u_0 = 1} & \boldsymbol{u_1 = 0} & u_2 = u_0 - (q_0 \cdot u_1) & u_3 = u_1 - (q_1 \cdot u_2) & \cdots\\
				\hline
					\boldsymbol{v_0 = 0} & \boldsymbol{v_1 = 1} & v_2 = v_0 - (q_0 \cdot v_1) & v_3 = v_1 - (q_1 \cdot v_2) & \cdots\\
				\hline
			\end{array} 
		\]

Les derniers éléments des deux dernières lignes sont les coefficients \(u\) et \(v\) qu'on cherche.
\end{frame}
\begin{frame}[allowframebreaks]{PPCM et PGCD}
En fait, à chaque étape, on a
		\(
			r_{n-2} = a_{n} = u_{n} \; a + v_{n} \; b
		\).
\end{frame}
\begin{frame}[allowframebreaks]{PPCM et PGCD}
Revenons à l'exemple précédent.
		On sait que \(PGCD(200,36) = 4\).
		On veut utiliser l'algorithme d'Euclide pour trouver les nombres \(u\) et \(v\) tels que 
		\(200 \, u + 36 \, v = 4\).
\end{frame}
\begin{frame}[allowframebreaks]{PPCM et PGCD (2)}
\begin{center}\begin{tabular}{lllll}
 & \(5\) & \(1\) & \(1\) & \(4\) \\ {}
\(200\) & \(36\) & \(20\) & \(16\) & \(4\) \\ {}
\(20\) & \(16\) & \(4\) & \(0\) &  \\ {}
\(1\) & \(0\) & \(1\) & \(-1\) & \(2\) \\ {}
\(0\) & \(1\) & \(-5\) & \(6\) & \(-11\) \\ {}
\end{tabular}\end{center}
\end{frame}
\begin{frame}[allowframebreaks]{PPCM et PGCD (3)}
Ainsi, on retrouve bien \(200 \cdot 2 - 36 \cdot 11 = 4\).
\end{frame}
\begin{frame}[allowframebreaks]{PPCM et PGCD}
Soient \(a\), \(b\) et \(c\) trois entiers relatifs tels que \(PGCD(a,c) = 1\).
                  Alors, \(a|bc \Rightarrow a | b\).
\end{frame}
\begin{frame}[allowframebreaks]{PPCM et PGCD}
On se propose de résoudre dans \(\mathbb{Z}\times\mathbb{Z}\) l'équation : 
                  \(35 x + 38 y = 1\).
                  
                  Commençons par dresser l'algorithme d'Euclide pour trouver \(PGCD(38,35)\) :
\end{frame}
\begin{frame}[allowframebreaks]{PPCM et PGCD (2)}
\begin{center}\begin{tabular}{lllll}
 & \(1\) & \(11\) & \(1\) & \(2\) \\ {}
\(38\) & \(35\) & \(3\) & \(2\) & \(1\) \\ {}
\(3\) & \(2\) & \(1\) & \(0\) &  \\ {}
\end{tabular}\end{center}
\end{frame}
\begin{frame}[allowframebreaks]{PPCM et PGCD (3)}
.
                  
                  On peut alors chercher deux entiers \(u\) et \(v\) tels que \(38 u + 35 v = 1\).
                                    \[
                        \begin{array}{|r@{\ =\ }l|r@{\ =\ }l|r@{\ =\ }l|}
                              \hline
                              \multicolumn{2}{|c|}{\text{Division}} &
                                     \multicolumn{2}{|c|}{\text{Ecriture numérique} } & 
                                           \multicolumn{2}{|c|}{\text{Écriture littérale} } \\
                              \hline
                              38  & 1 \cdot 35 + 3 & 3 & 38 - 1 \cdot 35       & r_0  & a   - 1 \cdot b    \\
                              \hline
                              35  & 11 \cdot 3 + 2 & 2 & 35  - 11 \cdot 3      & r_1  & b   - 11 \cdot r_0 \\
                              \hline
                              3   & 1 \cdot 2 + 1  & 1  & 3  - 1 \cdot 2       & r_2  & r_0 - 1 \cdot r_1  \\
                              \hline
                              2  & 2 \cdot 1  + 0  & \multicolumn{2}{|c|}{}      & \multicolumn{2}{|c|}{}  \\
                              \hline
                        \end{array}
                  \]
\end{frame}
\begin{frame}[allowframebreaks]{PPCM et PGCD (4)}
Ainsi, on a : 
                        \[\begin{aligned}
                              r_2 & = (a - b ) - (b - 11 \cdot r_0) \\
                                  & = a - b - b + 11 (a - b)        \\
                                  & = a - 2 b + 11 a - 11 b         \\
                                  & = 12 a - 13 b                   \;.
                        \end{aligned}\]

On aurait aussi pu utiliser l'algorithme d'Euclide étendu :
\end{frame}
\begin{frame}[allowframebreaks]{PPCM et PGCD (5)}
\begin{center}\begin{tabular}{lllll}
 & \(1\) & \(11\) & \(1\) & \(2\) \\ {}
\(38\) & \(35\) & \(3\) & \(2\) & \(1\) \\ {}
\(3\) & \(2\) & \(1\) & \(0\) &  \\ {}
\(1\) & \(0\) & \(1\) & \(-11\) & \(12\) \\ {}
\(0\) & \(1\) & \(-1\) & \(12\) & \(-13\) \\ {}
\end{tabular}\end{center}
\end{frame}
\begin{frame}[allowframebreaks]{PPCM et PGCD (6)}
On a donc une solution particulière : \(12 \cdot 38 - 13 \cdot 35 = 1\).
                  Si \(u\) et \(v\) sont deux solutions, on a : 
                        \[\begin{aligned}
                        &
                              \left\lbrace
                                    \begin{array}{l}
                                          38 \cdot  u + 35 \cdot v  = 1 \\
                                          38 \cdot 12 - 35 \cdot 13 = 1 \\
                                    \end{array}
                              \right.
                        \Rightarrow
                        &     
                                    38 (u - 12) + 35 (v + 13) = 0      \;.
                        \end{aligned}\]

Ainsi,
                              \[
                                    35 \, (v + 13) = 38 (12 - u) 
                              \]
\end{frame}
\begin{frame}[allowframebreaks]{PPCM et PGCD (7)}
\footnotesize
et donc \(38 | 35 \, (v + 13)\).
                        Mais comme \(PGCD(38,35) = 1\), d'après le Lemme de Gauss, \(38 | (v + 13)\).
                        D'où, il existe un entier \(k\) tel que \(v + 13 = 38 \, k\), ou encore \(v = 38\,k - 13\).
                        En portant cette valeur dans l'équation eq:Bezout ex 1, on a 
                              \(38 \, (12 - u) = 35 \cdot 38 \cdot k\) et donc 
                              \(u = - 35 \, k + 12\).
                        Réciproquement, pour tout \(k\in\mathbb{Z}\), le couple \((u ; v) = (12 - 35\,k ; 38\,k - 13)\) vérifie
                        \(38\,(12-35\,k) + 35\,(38\,k-13) = 456 - 455 = 1\).
                        Comme \(x = v\) et \(y = u\) (l'équation de départ est \(35 x + 38 y = 1\)),
                        l'ensemble des solutions est
                              \(\left\lbrace (38\,k-13 ; 12-35\,k) \mid k \in \mathbb{Z} \right\rbrace\).
\end{frame}
\begin{frame}[allowframebreaks]{PPCM et PGCD}
Deux entiers relatifs \(a\) et \(b\) sont premiers entre eux si et seulement si
                  il existe deux entiers relatifs \(u\) et \(v\) tels que \(a u + b v = 1\).
\end{frame}
\begin{frame}[allowframebreaks]{PPCM et PGCD}
\begin{enumerate}
\setcounter{enumi}{0}
                        \item Si un nombre premier \(p\) divise un produit \(ab\), alors \(p\) divise \(a\) ou \(p\) divise \(b\).
                        
\item Si \(a | n\), \(b | n\) et \(PGCD(a,b) = 1\), alors \(ab | n\).

\end{enumerate}
\end{frame}
\begin{frame}[allowframebreaks]{PPCM et PGCD}
Soient \(a\) et \(b\) deux entiers relatifs non nuls, \(d = PGCD(a,b)\) et \(c \in \mathbb{Z}\).
\end{frame}
\begin{frame}[allowframebreaks]{PPCM et PGCD (2)}
\begin{itemize}
\item L'équation \(a x + b y = c\) a des solutions dans \(\mathbb{Z}^2\) si et seulement si \(d\) divise \(c\).

\end{itemize}
\end{frame}
\begin{frame}[allowframebreaks]{PPCM et PGCD (3)}
\begin{itemize}
\item Si \(d | c\), on divise par \(d\) : \(a' x + b' y = c'\) avec \(a = d a'\), \(b = d b'\), \(c = d c'\) et \(PGCD(a',b') = 1\).
                        On cherche une solution particulière \((x_0 ; y_0)\) (par exemple en multipliant par \(c'\) des coefficients de Bézout de \(a'\) et \(b'\)),
                        puis le lemme de Gauss donne toutes les solutions :
                        \[
                              x = x_0 + b' k , \qquad y = y_0 - a' k , \qquad k \in \mathbb{Z} .
                        \]

\end{itemize}
\end{frame}
\begin{frame}[allowframebreaks]{PPCM et PGCD}
Soient \(n \geq 2\) et \(a\in\mathbb{Z}\).
\end{frame}
\begin{frame}[allowframebreaks]{PPCM et PGCD (2)}
\begin{enumerate}
\setcounter{enumi}{0}\item Il existe \(\overline{u}\) tel que \(\overline{a} \cdot \overline{u} = \overline{1}\) dans \(\mathbb{Z}/n\mathbb{Z}\)
                              si et seulement si \(PGCD(a,n) = 1\) (c'est le théorème de Bézout : \(au + nv = 1\)).
                              On peut alors résoudre \(\overline{a}\,x = \overline{b}\) en multipliant par \(\overline{u}\) : \(x = \overline{u}\,\overline{b}\).

\end{enumerate}
\end{frame}
\begin{frame}[allowframebreaks]{PPCM et PGCD (3)}
\begin{enumerate}
\setcounter{enumi}{1}\item Si \(p\) est premier, tout élément non nul de \(\mathbb{Z}/p\mathbb{Z}\) a un inverse, et
                              \(\overline{x} \cdot \overline{y} = \overline{0} \Rightarrow \left( \overline{x} = \overline{0} \text{ ou } \overline{y} = \overline{0} \right)\).
                              Ce n'est plus vrai si \(n\) n'est pas premier : dans \(\mathbb{Z}/6\mathbb{Z}\), \(\overline{2}\cdot\overline{3} = \overline{0}\).

\end{enumerate}
\end{frame}
\end{document}
